\documentclass[11pt]{article}
\usepackage[T1]{fontenc}
\usepackage{textcomp}
\usepackage[margin=0.75in]{geometry}
\usepackage{amsmath,amssymb,amsthm}
\usepackage{array,booktabs}
\usepackage{hyperref}
\setlength{\parindent}{0pt}
\newtheorem{theorem}{Theorem}
\theoremstyle{definition}
\newtheorem{definition}{Definition}
\title{Flow Proof Artifact: pythagoras}
\author{Generated by \texttt{flow doc proof}}
\date{Flow Proof Book}
\begin{document}
\maketitle
In a right triangle, the square on the hypotenuse equals the sum of squares on the legs.\\[0.5em]
\textbf{Source.} euclid — Elements, Book I, Proposition 47\\[0.5em]
\bigskip
\noindent{\large \textbf{Derived fact 1.} \textit{In a right triangle, the square on the hypotenuse equals the sum of squares on the legs} \hfill the Euclidean plane \cdot right triangle \cdot the Pythagorean relation holds}\par
\smallskip\noindent\textit{Coordinate: the Euclidean plane \cdot right triangle \cdot the Pythagorean relation holds}\par
\smallskip\noindent\textit{Source: euclid — Elements, Book I, Proposition 47}\par
\medskip
\noindent\textbf{Goal.} In a right triangle, the square on the hypotenuse equals the sum of squares on the legs\par
\noindent$\displaystyle c^{2} = a^{2} + b^{2}$\par
\medskip
\noindent
\begin{tabular}{@{} >{\raggedright\arraybackslash}p{0.06\textwidth} >{\raggedright\arraybackslash}p{0.47\textwidth} >{\raggedleft\arraybackslash}p{0.41\textwidth} @{}}
\textbf{\#} & \textbf{Proof} & \textbf{Mathematics} \\
\midrule
\textbf{1.} & We prove that the Pythagorean relation holds for right triangles in the Euclidean plane. & \\[0.45em]
\textbf{2.} & the square on the hypotenuse decomposes into the two squares on the legs, so c times c equals a times a plus b times b. Hence proven. & $\displaystyle c^{2} = a^{2} + b^{2}$ \\[0.45em]
\bottomrule
\end{tabular}
\label{thm:Geometry-right triangle-the_pythagorean_relation_holds}
\par\medskip\hrule
\end{document}
